\section{Jets and the three rows}\label{sec:jets}

We retain the prepared data of Proposition~\ref{prop:prepared-cycle}.
Thus $V$ is a smooth projective threefold, and
\[
 D=\sum_\alpha m_\alpha D_\alpha
\]
is an effective Cartier divisor with connected support,
$\gcd_\alpha m_\alpha=1$, and $D|_{D_\alpha}$ numerically trivial
for every $\alpha$. In particular, $D^2A=0$ for every divisor $A$.
The line bundle $\cO_D(D)$ is torsion. Near $D$ there is an
isomorphism
\begin{equation}\label{eq:canonical-p-power}
 \omega_V^{\otimes p^b}\simeq\cO_V(B_*),\qquad
 B_*=\sum_\alpha b_\alpha D_\alpha,
 \quad b_\alpha\in\Z,
\end{equation}
and $K_V^2D=0$. We shall use
$c_2(V)D\geq0$ from Proposition~\ref{prop:strong-semistability}
and the absence of scalar polar parts from
Corollary~\ref{cor:no-scalar-poles}.

The argument in this section concerns finite jets and their ranks.
Its two outputs are a trivialization of the canonical line on the
completion up to an integral divisor supported on $D$, and the equality
$c_2(V)D=0$ together with the limit $p/(p+1)$ for the normalized last
trivialization widths.  The width at period $rp^e$ is the largest
integer $s$ for which $\cO(rp^eD)$ is trivial on $sD$.
Its only passage to a formal limit will be an explicit inverse-limit
argument for finite-dimensional spaces of sections.

\subsection{Finite pages and duality}

Write $\widehat{\cO}_V$ for the completion along $D$, regarded as
a sheaf on $|D|$, and set
\[
 \cR=\widehat{\cO}_V(*D).
\]
Choose a finite cover by the restrictions to $|D|$ of affine open
subsets of $V$, refining it so that $D$ is principal and the bundles
under consideration are free on each member. All intersections of
these affine opens are affine, since $V$ is separated. For a vector
bundle $E$, let $C_i^\bullet(E)$ be the resulting \v{C}ech complex
with coefficients in the completion of $E(iD)$. These complexes
give an increasing filtration, indexed by $i\in\Z$, on the Laurent
\v{C}ech complex.

Here are the elementary completeness facts we need. On an affine
chart with local equation $f$ for $D$, adic completion is flat and
$f$ remains a non-zero-divisor. Thus the completed lattices embed
in their localization. A section of the localization on a
quasi-compact open lies in one lattice: finitely many local pole
bounds have a common upper bound, and the resulting lattice
sections glue. Consequently every Laurent cochain has a finite
upper weight. Each $C_i^\bullet(E)$ is complete and separated for
the filtration by $C_{i-s}^\bullet(E)$, $s\geq1$. Exactness of
completion for finite modules gives
\begin{equation}\label{eq:jet-quotient}
 \begin{gathered}
 C_i^\bullet(E)/C_{i-s}^\bullet(E)
 \quad\text{computes}\quad H^\bullet\bigl(V,P_i^s(E)\bigr),\\
 P_i^s(E)=\frac{E(iD)}{E((i-s)D)}.
\end{gathered}
\end{equation}
Indeed its terms are exactly the sections of this algebraic
quotient on the affine intersections. Thus $P_i^s(E)$ retains
the $s$ layers at weights $i-s+1,\ldots,i$; we also call $i$
the upper cut. We abbreviate
$E^!=E^\vee\otimes\omega_V$.

For $s\geq1$, consider a leading cochain, that is, a class in
$C_i^a(E)/C_{i-1}^a(E)$, where $a$ is its cohomological degree.
The cycle condition below asks whether this class has a representative
closed modulo $C_{i-s}^\bullet(E)$. The boundary condition asks whether
it is the leading term of a coboundary with a primitive in
$C_{i+s-1}^{a-1}(E)$. Thus the index $s$ records both the precision
of closedness and the allowed upper cut of a primitive. Define
\begin{align*}
 Z_s^{i,a}&=
 \{[x]:x\in C_i^a(E),\ \partial x\in C_{i-s}^{a+1}(E)\},\\
 U_s^{i,a}&=
 \{[\partial y]:y\in C_{i+s-1}^{a-1}(E),\
                         \partial y\in C_i^a(E)\}.
\end{align*}
Every exact cochain is closed, so $U_s^{i,a}\subseteq Z_s^{i,a}$.
Put
\[
 E_s^{i,a}(E)=Z_s^{i,a}/U_s^{i,a},\qquad
 g_s^a(i;E)=\dim_k E_s^{i,a}(E).
\]
The notation $E_s^{i,a}(E)$ denotes a page space, not a new bundle.

\begin{lemma}[Finite-page calculus]\label{lem:page-calculus}
The operation $x\mapsto\partial x$ induces maps
\[
 d_s:E_s^{i,a}(E)\longrightarrow E_s^{i-s,a+1}(E)
\]
whose cohomology is the next page. At page one,
$E_1^{i,a}(E)=H^a(V,P_i^1(E))$; every page has only rows
$a=0,1,2$. In particular, with
$\delta_s^a(i)=g_s^a(i;E)-g_{s+1}^a(i;E)$ for $a=0,2$,
\begin{align}
 g_s^0(i;E)
 &=\dim\im\bigl(H^0(P_i^s(E))\longrightarrow H^0(P_i^1(E))\bigr),
 \label{eq:bottom-endpoint}\\
 g_s^2(i;E)&=g_s^0(1-i;E^!),\label{eq:top-endpoint}\\
 g_{s+1}^1(i;E)
 &=g_s^1(i;E)-\delta_s^0(i+s)-\delta_s^2(i-s).
 \label{eq:middle-recurrence}
\end{align}
If $M>0$ and $\cO(MD)|_{sD}$ is trivial, the page-$s$ ranks are
periodic in $i$ modulo $M$, and
\begin{equation}\label{eq:periodic-euler}
 \sum_{i\bmod M}\bigl(g_s^0(i;E)-g_s^1(i;E)+g_s^2(i;E)\bigr)
 =\sum_{i\bmod M}\chi(P_i^1(E)).
\end{equation}
Finally, at each fixed weight $i$, the stable image in
Equation~\ref{eq:bottom-endpoint} is precisely the image of the
global sections of the completed bundle $\widehat{E(iD)}$ in
$H^0(P_i^1(E))$.
\end{lemma}

\begin{proof}
If two representatives of a leading cochain differ by
$y\in C_{i-1}^a$, and both have coboundary in $C_{i-s}$, their
coboundaries differ by a target boundary coming from
$C_{i-1}=C_{(i-s)+s-1}$. Adding an exact representative has no
effect either. This proves that $d_s$ is well defined, and
$d_s^2=0$ follows from $\partial^2=0$.

If $d_s[x]=0$, the leading term of $\partial x$ is represented
by $\partial y$ with $y\in C_{i-1}^a$ and
$\partial y\in C_{i-s}^{a+1}$. Replacing $x$ by $x-y$ makes its
coboundary lie in $C_{i-s-1}^{a+1}$. Thus the kernel is represented
by $Z_{s+1}^{i,a}$. The incoming image enlarges $U_s^{i,a}$ to
$U_{s+1}^{i,a}$: its representatives are exactly the leading
terms of $\partial z\in C_i^a$ with $z\in C_{i+s}^{a-1}$.
This proves the next-page assertion. Equation~\ref{eq:jet-quotient}
identifies page one. Since the quotients are supported on the
projective surface scheme $D$, their cohomology is zero outside
degrees $0,1,2$; the same is true of every subsequent page.

The differentials through a middle-row term are arranged as
\[
 \begin{tikzcd}[column sep=large]
 E_s^{i+s,0}(E) \arrow[r,"d_s"] &
 E_s^{i,1}(E) \arrow[r,"d_s"] & E_s^{i-s,2}(E).
 \end{tikzcd}
\]
Moving right lowers weight by $s$ and raises cohomological degree
by one; the composite is zero. There is no incoming differential
in degree zero. Its representatives are exactly the sections of
$P_i^s(E)$, proving Equation~\ref{eq:bottom-endpoint}.

For the other endpoint, consider the inclusion of the bottom
layer
\[
 H^2(P_i^1(E))\longrightarrow H^2(P_{i+s-1}^s(E)).
\]
There is no outgoing differential in degree two. More explicitly,
the kernel calculation above and vanishing of the degree-three
page imply successively that $Z_s^{i,2}=Z_1^{i,2}$.
A class represented by $x\in C_i^2$, with
$\partial x\in C_{i-1}^3$, dies under this inclusion exactly when
$x-\partial y\in C_{i-1}^2$ for some
$y\in C_{i+s-1}^1$. Its kernel is therefore
$U_s^{i,2}/U_1^{i,2}$, and the image has dimension $g_s^2(i;E)$.

The two-term locally free presentation of $P_j^s(E)$ gives
\[
 \mathcal E xt^1(P_j^s(E),\omega_V)
 =\frac{E^!((s-j)D)}{E^!(-jD)}=P_{s-j}^s(E^!),
\]
and all other sheaf Ext groups vanish. Serre duality on the
smooth projective threefold consequently identifies
$H^2(P_j^s(E))^\vee$ with $H^0(P_{s-j}^s(E^!))$
\cite[Chapter~III, Section~7]{hartshorne1977}.
Under this identification the dual of the displayed inclusion is
\[
 H^0(P_{1-i}^s(E^!))\longrightarrow H^0(P_{1-i}^1(E^!)).
\]
One can check the map, as well as its index, directly: the map of
the Cartier presentations is the identity on $E((i-1)D)$ and
the inclusion on their upper terms, so dualizing gives precisely
the indicated projection. This proves Equation~\ref{eq:top-endpoint}.
The ranks of the two differentials adjacent to the middle term
are $\delta_s^0(i+s)$ and $\delta_s^2(i-s)$, respectively.
Taking its cohomology proves Equation~\ref{eq:middle-recurrence}.

Triviality of $\cO(MD)$ on $sD$ identifies each quotient of width
$t\leq s$ with its translate by $M$, compatibly with its layer
maps. Equations~\ref{eq:bottom-endpoint} and
\ref{eq:top-endpoint} therefore give periodicity of both outer
rows at every such width. Page one is periodic in all degrees;
Equation~\ref{eq:middle-recurrence} then gives periodicity of
the middle row by induction. If
$e_t(i)=g_t^0(i)-g_t^1(i)+g_t^2(i)$, that equation yields
\[
 e_{t+1}(i)-e_t(i)=
 -\delta_t^0(i)-\delta_t^2(i)
 +\delta_t^0(i+t)+\delta_t^2(i-t).
\]
For $t<s$ all four functions are periodic modulo $M$.
Their sums cancel, proving Equation~\ref{eq:periodic-euler}
by induction from page one.

For the last assertion, set $A_s=H^0(P_i^s(E))$ and
\[
 B_s=\bigcap_{t\geq s}\im(A_t\longrightarrow A_s).
\]
Each $A_s$ is finite dimensional, so the intersection defining
$B_s$ is attained at a finite stage. Choosing one stage beyond
stabilization at both $s$ and $s+1$ shows that
$B_{s+1}\longrightarrow B_s$ is surjective. Every element of
$B_1$ therefore admits compatible lifts in all the $B_s$.
Their inverse limit is a global section of $\widehat{E(iD)}$,
since sections commute with inverse limits of sheaves. Conversely
a completed section restricts to the image of every $A_s$.
This proves the assertion directly.
\end{proof}

\subsection{Scalar jets and their last widths}

\begin{lemma}[The scalar jet rule]\label{lem:scalar-jets}
Let $u<i$ be integers and let
$\xi\in H^0(\cO(iD)/\cO(uD))$ be nonzero. There is an integer
$t$ with $u<t\leq i$ such that local lifts of $\xi$ have pole
order exactly $tm_\alpha$ at every $D_\alpha$. Those lifts
give a trivializing section of $\cO(tD)$ on $(t-u)D$.
In particular,
\begin{equation}\label{eq:scalar-endpoint}
 g_s^0(i;\cO_V)=
 \begin{cases}
 1,&\cO(iD)|_{sD}\simeq\cO_{sD},\\
 0,&\text{otherwise}.
 \end{cases}
\end{equation}
\end{lemma}

\begin{proof}
Choose local algebraic lifts $f_\lambda$ of $\xi$ on a finite
affine cover. Membership in $\cO(uD)$ is checked by valuations
at the prime divisors of the smooth, hence normal, variety $V$.
Thus some component detects the nonzero quotient section. The
maximum of the ratios
$-\ord_{D_\alpha}(f_\lambda)/m_\alpha$ that exceed $u$ is a
well-defined rational number $t\in(u,i]$: changing a lift adds
a section of $\cO(uD)$ and cannot alter such a valuation.
All lifts belong to $\cO(tD)$ in the sense of these rational
divisor bounds.

Choose a positive integer $n$ with $nt\in\Z$. Then
$f_\lambda^n\in\cO(ntD)$, and these sections agree on
$D_{\mathrm{red}}$. Indeed
$f_\lambda-f_\mu\in\cO(uD)$, and in the factorization of
$f_\lambda^n-f_\mu^n$ each term has one factor of pole ratio
at most $u$ and $n-1$ factors of pole ratio at most $t$.
After normalization in $\cO(ntD)$ its valuation at every
$D_\alpha$ is at least $(t-u)m_\alpha>0$.
They therefore define a section of
$\cO(ntD)|_{D_{\mathrm{red}}}$, nonzero on a component where
the maximum is attained.

The restriction of this line bundle to every integral component
is numerically trivial. A nonzero section of a numerically
trivial line bundle on a projective integral surface has no
zeros: a nonempty effective Cartier zero divisor has positive
intersection with an ample divisor. It is consequently nowhere
zero on that component. At an intersection with another
component its common value is nonzero, so its restriction to
that component is also nonzero and nowhere zero. Connectedness
of $|D|$ propagates this conclusion to all components.

It follows that $\ord_{D_\alpha}(f_\lambda)=-tm_\alpha$
for every $\alpha$. In particular every $tm_\alpha$ is an
integer. Primitivity of the cycle, or an integral linear
combination of the $m_\alpha$ equal to one, gives $t\in\Z$.
The original lifts are now sections of the honest line bundle
$\cO(tD)$, are units in it along $D$, and agree modulo
$\cO(uD)$. Since
\[
 \frac{\cO(tD)}{\cO(uD)}
 =\cO(tD)|_{(t-u)D},
\]
they give the asserted trivialization.

A jet of width $s$ at upper weight $i$ has nonzero image in
the leading layer exactly when the uniform ratio just proved
is $t=i$. This proves existence in
Equation~\ref{eq:scalar-endpoint}, in both directions.
To see that the image dimension is one, take two such jets.
Their normalized unit sections on $D_{\mathrm{red}}$ are
proportional: a connected reduced projective scheme over $k$
has only constant global functions. Subtract the appropriate
multiple. If its leading image were still nonzero, the preceding
argument would again give a unit on the reduction with ratio
$i$, contradicting the cancellation. Thus the two leading
images are proportional.
\end{proof}

The organizing viewpoint of tracking torsion and triviality on infinitesimal
thickenings was informed by Totaro's treatment of these obstructions
\cite[Lemma~4.1 and Section~6]{totaro2007}. The elementary width properties
are proved below; the later three-row rank calculation and Cartier descent
are not inputs from that work.

Let $r$ be the order of $\cO_D(D)$. For $e\geq0$ put
\begin{equation}\label{eq:width-definition}
 M_e=rp^e,\qquad
 d_e=\max\{s\geq1:\cO(M_eD)|_{sD}\simeq\cO_{sD}\}.
\end{equation}
These maxima exist and satisfy
\begin{equation}\label{eq:width-bounds}
 1\leq d_e<M_e,\qquad d_{e+1}\geq p d_e,
 \qquad \frac{d_e}{M_e}\nearrow a_*\in(0,1].
\end{equation}
In fact a trivialization on $M_eD$ would give, by
Lemma~\ref{lem:scalar-jets}, a nonzero section of
$\cO(M_eD)/\cO$, contrary to
Corollary~\ref{cor:no-scalar-poles}. Triviality at width one
follows from the definition of $r$. Taking $p$th powers of
local trivializing sections multiplies the gluing precision
by $p$ and gives $d_{e+1}\geq pd_e$. Division by
$M_{e+1}=pM_e$ proves monotonicity and the positive limit.
If $r=1$, already $1\leq d_0<M_0=1$ is a contradiction;
we may therefore continue with $r>1$.

We will repeatedly use the following precise consequence of
nilpotence. If a line bundle $L$ on a finite thickening $T$
of $D$ is trivial on $D$, then its class has $p$-power order.
Indeed its transition functions can be chosen in $1+J$, where
$J$ is the nilpotent ideal of $D$ in $T$. For $p^N$ at least
a nilpotence exponent, $(1+x)^{p^N}=1$ for $x\in J$.
Thus $L^{\otimes p^N}$ is trivial. Multiplication by an
integer prime to $p$ is injective on such classes, by
B\'ezout's identity, as well as surjective on their cyclic
subgroups. It follows that
\begin{equation}\label{eq:scalar-last-width}
 \text{the last width at }i=jM_e\text{ is }d_e
 \quad\text{if }p\nmid j.
\end{equation}
This includes negative $j$. If $r\nmid i$, no scalar leading
jet occurs even at width one. Every nonzero multiple of $r$
has a unique expression $jM_e$ with $p\nmid j$; at $i=0$
the scalar entry persists at every width. These observations
classify all scalar outer-row entries and drops.

Finally Hirzebruch--Riemann--Roch on the smooth threefold gives
\begin{equation}\label{eq:scalar-euler}
 \chi(P_i^1(\cO_V))
 =\chi(\cO(iD))-\chi(\cO((i-1)D))
 =\frac{Dc_2(V)}{12}=:\gamma\geq0.
\end{equation}
Indeed the terms involving $D^3$, $D^2K_V$, and $DK_V^2$
vanish, leaving the displayed linear term in the threefold
Riemann--Roch polynomial \cite[Section~15]{fulton1998}.

\subsection{Canonical jets and a formal canonical root}

For $E=\cO_V$, Equation~\ref{eq:top-endpoint} computes the top-row
rank at weight $i$ from leading jets of $\omega_V$ at weight $1-i$.
We need to describe that row as
explicitly as the scalar bottom row. We first show that canonical jets
survive to arbitrarily large widths. Their $p^b$th powers then let the
scalar rule identify $\widehat{\omega}_V$ with the line of an integral
divisor supported on $D$.

\begin{lemma}\label{lem:forced-canonical-jets}
There are nonzero leading jets of $\omega_V$ of arbitrarily
large width, allowing their upper weight to vary.
\end{lemma}

\begin{proof}
Suppose instead that $g_s^0(l;\omega_V)=0$ for every $l$
and every $s\geq S$, for some $S$. By
Equation~\ref{eq:top-endpoint}, the top row for $E=\cO_V$
then vanishes at all these widths. Take $e$ sufficiently
large that $s_e=d_{e-1}+1\geq S$. At page $s_e$ the bottom
row is one precisely at multiples of $M_e$, by
Equation~\ref{eq:scalar-last-width} and
$s_e\leq d_e$. This page is periodic modulo $M_e$.
Equations~\ref{eq:periodic-euler} and
\ref{eq:scalar-euler} give
\[
 \sum_{i\bmod M_e}g_{s_e}^1(i;\cO_V)
 =1-M_e\gamma\leq1.
\]
At differential length $d_e$, the bottom entries at $jM_e$,
$p\nmid j$, each drop by one. By
Equation~\ref{eq:middle-recurrence} these drops consume one
middle-row dimension at $-d_e+jM_e$. Middle-row dimensions
can only decrease as the page increases. Periodicity and
the preceding upper bound therefore force the initial middle
row to consist exactly of one entry at
$-d_e\pmod{M_e}$.

Between pages $s_e$ and $s_{e+1}=d_e+1$ these are the only
drops: there is no top row, and all scalar entries with a
different last-width index either disappeared earlier or have
last width at least $d_{e+1}$. Modulo $M_{e+1}=pM_e$, the
drops remove all the $p$ lifts of the middle entry except its
unshifted lift. Thus at page $s_{e+1}$ the middle row consists
of one entry at $-d_e\pmod{M_{e+1}}$. Applying the same
argument with $e+1$ says that it must instead lie at
$-d_{e+1}\pmod{M_{e+1}}$. But
$0<d_e<d_{e+1}<M_{e+1}$, so these residues are different.
This contradiction proves the lemma.
\end{proof}

\begin{proposition}[The canonical formal root]\label{prop:canonical-formal-root}
Every coefficient of $B_*$ in
Equation~\ref{eq:canonical-p-power} is divisible by $p^b$.
Writing $B_*=p^bB_0$ with $B_0$ integral, there is an
isomorphism on the completion along $D$,
\begin{equation}\label{eq:canonical-formal-root}
 \widehat{\omega}_V\simeq\widehat{\cO}_V(B_0).
\end{equation}
\end{proposition}

\begin{proof}
Put $P=p^b$, and choose a sequence of nonzero leading
canonical jets of widths $s\to\infty$, with upper weights
$l$ that may vary. Raising local representatives to the
$P$th power and using Equation~\ref{eq:canonical-p-power}
gives a scalar jet with bounds
\[
 A=PlD+B_*,\qquad B=P(l-s)D+B_*.
\]
The lower bound is multiplied by $P$ exactly because $P$
is a power of the characteristic. Set
$\beta_\alpha=b_\alpha/m_\alpha$ and choose an integer
$c_*\geq\max_\alpha\beta_\alpha$. Enlarging both bounds
gives a section of
\[
 \frac{\cO((Pl+c_*)D)}{\cO((P(l-s)+c_*)D)}.
\]
This section is nonzero for all sufficiently large $s$,
uniformly in $l$. To check this, at a component detecting
the original leading layer, a local coefficient $f$ in
an ordinary frame of $\omega_V$ has
$-\ord_{D_\alpha}(f)/m_\alpha>l-1$. Its powered scalar
coefficient has pole ratio
\[
 -P\frac{\ord_{D_\alpha}(f)}{m_\alpha}
       +\beta_\alpha>P(l-1)+\min_\alpha\beta_\alpha.
\]
The new lower bound is $P(l-s)+c_*$, which is smaller for
large $s$.

Lemma~\ref{lem:scalar-jets} now shows that all these powered
coefficients have one integral pole ratio $t$. The preceding
lower estimate and the upper bound give a constant $C$, independent
of $s,l$, such that
\begin{equation}\label{eq:canonical-ratio-bound}
 \begin{gathered}
 |t-Pl|\leq C,\\
 \cO(tD)\text{ is trivial to a width at least }Ps-C.
 \end{gathered}
\end{equation}
For each fixed $J$, widths greater than $d_{J-1}$ can occur
only at multiples of $M_J$, by the scalar classification;
zero is allowed as a multiple. Since the widths in
Equation~\ref{eq:canonical-ratio-bound} tend to infinity,
we can assign integers $J\to\infty$ along this sequence
such that $M_J\mid t$.

Fix a component $D_\alpha$. If $f$ is the original local
coefficient in a frame of $\omega_V$, the local generator
of $\cO(B_*)$ has order $-b_\alpha$. Consequently
\begin{equation}\label{eq:canonical-order-identity}
 b_\alpha-tm_\alpha=P\ord_{D_\alpha}(f).
\end{equation}
For large $J$, the integer $t$ is divisible by $P$.
Equation~\ref{eq:canonical-order-identity} therefore shows
that $P\mid b_\alpha$ for every $\alpha$. This proves
$B_*=PB_0$ with $B_0$ integral.

For the same sequence put $h_0=t/P$. The order identity
says that the original local representatives are sections of
\[
 L_{h_0}=\omega_V(-B_0+h_0D)
\]
near $D$: their coefficient orders are exactly those
required by this line bundle. They have no additional poles
near $D$, since the original representatives were sections
of $\omega_V(lD)$ and membership in these divisor lattices
is checked on the normal variety in codimension one.
Moreover $L_{h_0}^{\otimes P}\simeq\cO(tD)$, and the
$P$th powers of these sections are units along all of $D$
by Lemma~\ref{lem:scalar-jets}. The original sections are
therefore themselves units along $D$.
They agree modulo the original lower lattice
$\omega_V((l-s)D)$. Hence they trivialize $L_{h_0}$ on
the effective thickness
\[
 T=(h_0-l+s)D-B_0.
\]
By Equation~\ref{eq:canonical-ratio-bound}, $|h_0-l|$ is
bounded independently of the sequence. Thus $T$ contains
$nD$ for integers $n\to\infty$. Also $M_J\mid t$
implies $M_{J-b}\mid h_0$ once $J\geq b$.
The line bundle $\cO(h_0D)$ is therefore trivial on
$d_{J-b}D$, and $d_{J-b}\to\infty$ by
Equation~\ref{eq:width-bounds}. Removing this twist shows
that $L=\omega_V(-B_0)$ is trivial on arbitrarily large
finite thickenings $nD$.

We finally choose these trivializations compatibly. Apply
the stabilized-image argument in the proof of
Lemma~\ref{lem:page-calculus} to the spaces
$H^0(nD,L|_{nD})$, also including restriction to
$D_{\mathrm{red}}$. The latter line bundle is trivial,
and its section space is one dimensional because
$D_{\mathrm{red}}$ is connected, reduced, and projective.
The image from every sufficiently large thickening contains
a unit, hence is that full one-dimensional space. The
surjective maps between stabilized images therefore lift a
chosen nonzero section on $D_{\mathrm{red}}$ compatibly to
all thickenings. Each lift is a unit, as may be checked after
reduction, and their inverse limit trivializes $\widehat L$.
This gives Equation~\ref{eq:canonical-formal-root}.
\end{proof}

\subsection{The two endpoint residues and the limiting width}

Write $B_0=\sum_\alpha b_{0,\alpha}D_\alpha$, and now use
$h$ for the integer offset
\begin{equation}\label{eq:canonical-offset}
 h=-\left\lfloor\min_\alpha
                   \frac{b_{0,\alpha}}{m_\alpha}\right\rfloor,
 \qquad q=1-h.
\end{equation}
The formal isomorphism of
Proposition~\ref{prop:canonical-formal-root} identifies the
finite canonical jet lattices with scalar lattices having
upper and lower bounds $lD+B_0$ and $(l-s)D+B_0$.
All these finite quotients are algebraic sheaves supported
on finite thickenings of $D$.
Indeed the localized $\widehat{\cO}_V$-linear isomorphism
respects every nested divisor lattice. Each width-$s$
quotient is killed by $\cI_D^s$, where $\cI_D=\cO_V(-D)$;
since $\widehat{\cO}_V/\cI_D^s\widehat{\cO}_V=\cO_{sD}$,
the induced map is an algebraic $\cO_{sD}$-linear map,
compatible with the leading-layer projections.

We first give the exact canonical counterpart of the scalar
rule. For all sufficiently large $s$, a canonical jet at
weight $l$ with nonzero leading image has uniform integral
scalar pole ratio
\begin{equation}\label{eq:canonical-scalar-ratio}
 t=l-h.
\end{equation}
To prove this, enlarge the two bounds to uniform multiples
of $D$ as in the preceding proof. A component detecting the
leading image has ratio within a fixed distance of $l$, so
for sufficiently large $s$ the enlarged quotient still
detects it. Lemma~\ref{lem:scalar-jets} gives a uniform
integral ratio $t$. Set
$\beta_0=\min_\alpha b_{0,\alpha}/m_\alpha$. Fitting
under every upper bound, but not under the preceding layer,
is exactly the pair of inequalities
\[
 l-1+\beta_0<t\leq l+\beta_0.
\]
Its unique integral solution is
$t=l+\lfloor\beta_0\rfloor=l-h$.

The normalized unit sections of $\cO(tD)$ glue on the
irregular thickness
\begin{equation}\label{eq:irregular-thickness}
 T_s=(s-h)D-B_0,
\end{equation}
because this is the difference between $tD$ and the lower
bound $(l-s)D+B_0$. Conversely a trivialization of
$\cO(tD)$ on $T_s$, with $t=l-h$, includes in the canonical
quotient and has nonzero leading image. Indeed $tD\leq
lD+B_0$, whereas the inequalities above prevent inclusion
in the preceding upper lattice. The leading image has
dimension one: normalized unit sections are proportional
on $D_{\mathrm{red}}$, and their difference cannot retain
the same leading ratio after that reduction cancels, by the
scalar rule. We have proved
\begin{equation}\label{eq:canonical-endpoint-rule}
 g_s^0(l;\omega_V)=
 \begin{cases}
 1,&\cO((l-h)D)|_{T_s}\simeq\cO_{T_s},\\
 0,&\text{otherwise},
 \end{cases}
 \qquad(s\text{ sufficiently large}).
\end{equation}

The comparison with ordinary widths is uniform and quite
explicit. Put
\[
 C_0=h+\left\lceil\max_\alpha
                    \frac{b_{0,\alpha}}{m_\alpha}\right\rceil
 \geq0.
\]
Then, coefficient by coefficient,
\begin{equation}\label{eq:thickness-comparison}
 (s-C_0)D\leq T_s\leq sD.
\end{equation}
For $e$ sufficiently large, let $d'_e$ be the last width
$s$ for which $\cO(M_eD)|_{T_s}$ is trivial. The thicknesses
are nested as $s$ increases. Equation~\ref{eq:thickness-comparison}
and the definition of $d_e$ give
\begin{equation}\label{eq:canonical-width-comparison}
 d_e\leq d'_e\leq d_e+C_0.
\end{equation}
In particular this last width is finite and exists in the
range where Equation~\ref{eq:canonical-endpoint-rule}
applies. Once $T_s$ contains $D$, the nilpotent-kernel
argument used above shows that the last width is the same
for $t=jM_e$ whenever $p\nmid j$. Also such a trivialization
requires $r\mid t$; for $t\ne0$ the unique expression
$t=jM_e$, $p\nmid j$, classifies its last width. The value
$t=0$ persists forever. To check that the large-$e$ rules
exhaust the entries at large width, observe from
Equation~\ref{eq:thickness-comparison} that an entry at
width $s$ gives an ordinary scalar trivialization to width
$s-C_0$. For each fixed $J$, taking $s-C_0>d_{J-1}$
forces $M_J\mid t$. Thus every nonzero index occurring at
sufficiently large width belongs to the stated range of
last-width indices.

\begin{proposition}[The limiting scalar widths]\label{prop:scalar-widths}
The scalar widths satisfy
\begin{equation}\label{eq:limiting-width}
 a_*=\lim_{e\to\infty}\frac{d_e}{M_e}
 =\frac{p}{p+1},\qquad Dc_2(V)=0.
\end{equation}
\end{proposition}

\begin{proof}
For all sufficiently large $e$ set
\[
 s_e=\max(d_{e-1},d'_{e-1})+1.
\]
Equations~\ref{eq:width-bounds} and
\ref{eq:canonical-width-comparison} show that these are
strictly increasing and
$s_e\leq\min(d_e,d'_e)$. Work with the three rows for
$E=\cO_V$. At page $s_e$ all three rows are periodic
modulo $M_e$, and each outer row has a single entry per period. The bottom
entry is at $0$, by Equation~\ref{eq:scalar-last-width}.
For the top entry, Equations~\ref{eq:top-endpoint} and
\ref{eq:canonical-scalar-ratio} associate to a top weight
$i$ the scalar ratio $t=1-i-h=q-i$. Its entry is therefore
at $q\pmod{M_e}$.

As the page runs from $s_e$ through $s_{e+1}$, the only
outer drops are those with last-width index $e$.
The bottom drops, of length $d_e$, consume middle-row
dimensions at
\[
 -d_e+jM_e,\qquad p\nmid j.
\]
The top drops, of length $d'_e$, consume middle-row
dimensions at
\[
 q+d'_e+jM_e,\qquad p\nmid j;
\]
here replacing $j$ by $-j$ accounts for the sign in the
top ratio. Put
\[
 A_e=-d_e,\qquad B_e=q+d'_e.
\]
The two lists, including their multiplicities when they
meet, are thus $A_e+jM_e$ and $B_e+jM_e$, with $p\nmid j$.

We spell out the rank count also when the two centers
coincide modulo $M_e$. Fix a residue $a\pmod{M_e}$ and
let $m(a)\in\{0,1,2\}$ be its multiplicity in the multiset
$\{A_e,B_e\}$. Its initial middle-row rank $n(a)$ is
constant on all $p$ lifts modulo $pM_e$. Each of the
$m(a)$ drop lists omits exactly one lift, namely its
unshifted center modulo $pM_e$. At most two lifts are
omitted.

\begin{samepage}
For example, when $p=5$ and $m(a)=2$, two
distinct omitted lifts can be pictured as follows:
\[
\begin{array}{c|ccccc}
 & a&a+M_e&a+2M_e&a+3M_e&a+4M_e\\
 \text{first list}&\circ&\bullet&\bullet&\bullet&\bullet\\
 \text{second list}&\bullet&\bullet&\circ&\bullet&\bullet
\end{array}
\]
A filled dot consumes one rank; a circle marks that list's
unshifted lift. The omitted lifts are illustrative. If they
coincide, their column is omitted twice and every other
column receives both drops. Multiplicities are always added.
\end{samepage}

In general, since $p>2$, some lift receives all $m(a)$ drops.
Nonnegativity of its final rank gives $n(a)\geq m(a)$.
Summing over $a$, and then using
Equations~\ref{eq:periodic-euler} and
\ref{eq:scalar-euler}, gives
\[
 2\leq\sum_{a\bmod M_e}n(a)
 =2-M_e\gamma\leq2.
\]
Consequently $\gamma=0$, so $Dc_2(V)=0$, and
$n(a)=m(a)$ for every residue.

After subtracting the drop lists, the final rank at each
lift is exactly the number of lists omitting that lift.
Thus at page $s_{e+1}$ the middle row is described by the
same two-element multiset $\{A_e,B_e\}$, now modulo
$M_{e+1}=pM_e$. Applying the just-proved description at
index $e+1$ gives the multiset equality
\begin{equation}\label{eq:residue-swap}
 \{-d_e,q+d'_e\}
 =\{-d_{e+1},q+d'_{e+1}\}\pmod{M_{e+1}}.
\end{equation}
The entries $-d_e$ and $-d_{e+1}$ are distinct modulo
$M_{e+1}$ because $0<d_e<d_{e+1}<M_{e+1}$.
Therefore the matching in Equation~\ref{eq:residue-swap}
must be the swapped matching, even if a multiset initially
has repeated entries. In particular
\[
 d_{e+1}+q+d'_e\equiv0\pmod{M_{e+1}}.
\]
For large $e$ its left side is positive and, by
$d_{e+1}<M_{e+1}$ and
$d'_e\leq d_e+C_0<M_e+C_0$, is less than
$2M_{e+1}$. Hence
\begin{equation}\label{eq:width-swap-equality}
 d_{e+1}+q+d'_e=M_{e+1}.
\end{equation}
Divide by $M_{e+1}=pM_e$ and let $e\to\infty$.
Equation~\ref{eq:canonical-width-comparison} gives
$d'_e/M_{e+1}\to a_*/p$, while
$d_{e+1}/M_{e+1}\to a_*$. Thus
$a_*+a_*/p=1$, which proves
Equation~\ref{eq:limiting-width}.
\end{proof}
